\documentclass[a4paper,10pt]{article}
\usepackage[utf8]{inputenc}
\usepackage[T1]{fontenc}
\usepackage{lmodern}
\usepackage[margin=1.4cm, top=1.2cm, bottom=1.2cm]{geometry}
\usepackage{hyperref}
\usepackage{xcolor}
\usepackage{enumitem}
\usepackage{graphicx}
\usepackage{titlesec}

% Minimal clean styling
\hypersetup{colorlinks=true, urlcolor=blue!70!black, linkcolor=blue!70!black}
\pagestyle{empty}
\setlength{\parindent}{0pt}

% Section formatting
\titleformat{\section}{\large\bfseries\color{black}}{}{0em}{}[\titlerule]
\titlespacing*{\section}{0pt}{6pt}{4pt}

% Colors
\definecolor{accent}{RGB}{20, 50, 120}

% Role command
\newcommand{\role}[4]{%
  \vspace{2pt}\noindent
  \begin{minipage}[t]{0.70\textwidth}\raggedright\textbf{#1} \textbullet\ \textit{#3}\end{minipage}%
  \hfill
  \begin{minipage}[t]{0.28\textwidth}\raggedleft\footnotesize\color{gray}#2 \textbullet\ #4\end{minipage}%
  \vspace{1pt}%
}
\setlist[itemize]{leftmargin=1.2em, itemsep=1.5pt, topsep=2pt, parsep=0pt}

\begin{document}

% ================= HEADER =================
\begin{minipage}[c]{0.80\textwidth}
  {\Huge\bfseries Prabhavit Gupta}\\[3pt]
  {\color{accent}\large AI Engineer $\vert$ AI Generativa, Architetture Agentiche, Azure AI Foundry $\&$ RAG}\\[6pt]
  {\footnotesize +39 351 610 7807 ~$\vert$~
  \href{mailto:prabhavitg@gmail.com}{prabhavitg@gmail.com} ~$\vert$~
  \href{https://www.linkedin.com/in/prabhavit/}{linkedin.com/in/prabhavit} ~$\vert$~
  \href{https://github.com/prab-gupta}{github.com/prab-gupta} ~$\vert$~
  Venezia, Italia (Residente Permanente UE $\vert$ Full Remote $\vert$ Disponibilità Immediata)}
\end{minipage}\hfill

% ================= SUMMARY =================
\section{Summary}
AI Engineer con 3 anni di esperienza commerciale nella progettazione, prototipazione e rilascio in produzione di soluzioni end-to-end di intelligenza artificiale generativa e architetture agentiche. Ampia competenza nell'integrazione di Large Language Models (LLM), sistemi RAG avanzati e servizi cloud-native (Azure AI, Azure AI Foundry, AWS Bedrock, GCP). Specializzato nella costruzione di flussi multi-agente autonomi con tool-calling, API integration e orchestrazione affidabile di chiamate asincrone. Certificato Microsoft nella creazione di agenti AI (Copilot Studio Agent Academy e Introduction to Building AI Agents). Residente permanente a Venezia, immediatamente disponibile per ingaggi full remote in Italia.

% ================= EXPERIENCE =================
\section{Experience}

\role{\href{https://flowspace.co}{FlowSpace.co}}{Mar 2024 - Aug 2026}{Lead AI Engineer $\&$ Agentic Systems Architect (Freelance)}{Londra, UK (remoto)}
\begin{itemize}
  \item Architettato e implementato l'intero ciclo di vita di sistemi AI generativi e agentici per oltre 4.500 utenti enterprise attivi, dalla prototipazione rapida al monitoraggio in produzione.
  \item Progettato workflow multi-agente autonomi dotati di tool calling, recupero RAG contestuale su Vector DB e fallback deterministici su API esterne.
  \item Integrato modelli linguistici di frontiera tramite API REST e SDK cloud (Azure OpenAI, Anthropic, Bedrock), ottimizzando costi di inferenza e latenza operativa (<1.8s).
  \item Implementato protocolli rigorosi di valutazione (evals), safety checks, tracciamento delle chiamate e monitoraggio continuo per eliminare allucinazioni.
  \item Collaborato con team di ingegneria e stakeholder aziendali per integrare soluzioni AI all'interno di architetture a microservizi scalabili.
  \item \textit{Tecnologie: Python, Azure AI, Azure OpenAI, LLM, RAG, LangGraph, FastAPI, Docker, PostgreSQL, MLOps, Git}
\end{itemize}

\role{Urlink}{Apr 2025 - Jul 2025}{AI $\&$ Software Engineer (Freelance)}{Barcellona, Spagna (remoto)}
\begin{itemize}
  \item Sviluppato pipeline per l'ingestion, embedding e retrieval semantico di flussi informativi non strutturati.
  \item Realizzato microservizi API in Python per integrare modelli generativi nei prodotti core con validazione schemi rigorosa.
  \item Ottimizzato le prestazioni di elaborazione dati e gestito la transizione da prototipo a deployment cloud-native.
  \item \textit{Tecnologie: Python, OpenAI/Azure APIs, Vector DB, PostgreSQL, Docker, REST APIs, Git}
\end{itemize}

\role{Consulente Tecnico Indipendente}{Sep 2023 - Mar 2024}{AI Solutions Engineer}{Venezia, Italia}
\begin{itemize}
  \item Progettato e realizzato soluzioni di automazione AI, assistenti conversazionali intelligenti e pipeline RAG per PMI italiane ed europee.
  \item Guidato le aziende clienti nell'adozione di architetture agentiche per l'automazione di workflow documentali e gestionali complessi.
  \item Definito architetture cloud sicure, gestito il deployment e fornito documentazione operativa e runbook tecnici.
  \item \textit{Tecnologie: Python, Azure AI Services, LangChain, LlamaIndex, Vector DB, FastAPI, Linux, Git}
\end{itemize}

\role{\href{https://tractionboard.co}{Tractionboard.co}}{Aug 2023 - Oct 2023}{Full-Stack Automation Developer (Stage)}{Utrecht, Paesi Bassi (remoto)}
\begin{itemize}
  \item Sviluppato webhook e flussi di integrazione API per piattaforme SaaS di monitoraggio analytics.
  \item \textit{Tecnologie: Python, JavaScript, REST APIs, Git}
\end{itemize}

% ================= EDUCATION =================
\section{Education}
\role{Università Ca' Foscari di Venezia}{Ott 2020 - Giu 2024}{Laurea Triennale in Filosofia, Studi Internazionali ed Economia, Minor in Imprenditorialità}{Venezia, Italia}

% ================= CERTIFICATIONS =================
\section{Certifications}
Microsoft: Introduction to building AI agents \textbullet\ Microsoft Copilot Studio Agent Academy \textbullet\ Microsoft Power Automate \textbullet\ Tray.AI Build Practitioner (Sep 2024) \textbullet\ Tray.AI Build Foundations

% ================= ADDITIONAL INFORMATION =================
\section{Additional Information}
\begin{itemize}
  \item \textbf{Inquadramento:} Residente permanente UE a Venezia, Italia. Disponibile per contratti CCNL o collaborazione con Partita IVA attiva.
  \item \textbf{Lingue:} Italiano (lavorativo), Inglese (fluente C1 professionale), Hindi (madrelingua).
  \item \textbf{Disponibilità:} Immediata (0 giorni di preavviso) per posizioni full remote.
\end{itemize}

% ================= CORE SKILLS =================
\section{Core Skills}
\begin{itemize}
  \item \textbf{AI Generativa $\&$ Sistemi Agentici:} Azure AI Foundry, Azure AI Services, Azure OpenAI, LLM, Agenti Multi-Agente, Tool Calling.
  \item \textbf{RAG $\&$ Vector Databases:} Retrieval-Augmented Generation, Chunking, Hybrid Search, Qdrant, Pinecone, pgvector.
  \item \textbf{Linguaggi $\&$ Integrazione API:} Python 3.12, FastAPI, Pydantic, REST API, Webhooks, JSON Schema, SQL.
  \item \textbf{Cloud $\&$ DevOps:} Architetture Cloud-Native, Docker, Kubernetes, Azure, AWS, CI/CD, Git, Linux.
  \item \textbf{Ciclo di Vita AI:} Prototipazione rapida, Testing, Evals, Deployment in produzione, Monitoraggio prestazioni e costi.
\end{itemize}

\end{document}
